\section{Perturbations and spectral rotation identities}\label{sec:rotation}

We now determine the overlap in each spectral subspace from the total
overlap.  Two ingredients enter: a small perturbation makes subsequential
overlap laws satisfy the Ghirlanda--Guerra identities, and differentiation
under Haar rotations then gives equations for the individual spectral
overlaps.  We prove these statements for a general deterministic spin
prior, since the same identities will also be needed on magnetization
slices.

\subsection{Enrichment and perturbations}

Throughout this section, the spectral alphabet
\(\lambda_1,\ldots,\lambda_m\) is fixed and bounded.
Write \(I_{a,N}\) for the spectral coordinates assigned to group \(a\),
and suppose that \(\rho_{a,N}=|I_{a,N}|/N\) tends to \(\rho_a>0\).
The prior \(\nu_N\) is any deterministic probability measure on
\(\{-1,1\}^N\), independent of the Haar matrix \(U\) and held fixed
as the parameters in this section vary.
A finite nonzero positive measure is also allowed after normalization: its
logarithmic mass is a deterministic addition to the log partition
function.  In particular, neither permutation symmetry of the prior nor
full support on the cube is assumed.

Fix a step partition with interval lengths \(w_i>0\),
\(0\le i<k\), and cascade exponents \(\zeta_i=\sum_{l<i}w_l\).
Let \((v_\alpha)\) be the associated cascade weights.  Write
\(c(\alpha,\alpha')\in\{0,\ldots,k-1\}\) for the common level of two
leaves, including \(k-1\) for equal leaves, and define
\[
 T_{\alpha\alpha'}=\frac{c(\alpha,\alpha')}{k}.
\]
This is a Gram kernel, with diagonal \((k-1)/k\).
For \(0\le h_0\le\cdots\le h_{k-1}\), let
\(z_i(\alpha)\), \(1\le i\le N\), be independent hierarchical Gaussian
fields with covariance
\[
 \E z_i(\alpha)z_j(\alpha')
   =\ind_{\{i=j\}}h_{c(\alpha,\alpha')}.
\]
They are independent of \(U\) and of the cascade weights.  The enriched
Hamiltonian, before the small perturbations, is
\[
 \frac t2\sigma^{\mathsf T}U^{\mathsf T}\Lambda_NU\sigma
 +\sum_{i=1}^N z_i(\alpha)\sigma_i-\frac N2h_{k-1},
 \qquad 0\le t\le1.
\]
For the uniform cube prior and \(t=0\), its expected pressure is
\(\ell(h)\), by factorization in the backward recursion.  The
unenriched model is the case \(k=1\), \(h_0=0\).

Enumerate the monomial kernels
\(\kappa_j\) in \((A^1,\ldots,A^m,T)\), with nonnegative integer
degrees and degree at most \(C(1+j)\).  The constant kernel may be
omitted.  These are Gram kernels, bounded by one in absolute value.
Conditional on \(U\), take independent centered Gaussian fields \(Y_j\)
with covariance \(\kappa_j\), independent of the cascade weights and the
linear enrichment.
One concrete realization uses coefficient tensors applied to
\[
 \bigotimes_{a=1}^m
 \left(\frac{(U\sigma)_{I_{a,N}}}{\sqrt N}\right)^{\otimes r_a}
\]
for the spin part of a monomial.  For a factor \(T^r\), use independent
copies of that tensor at the tree levels with variances given by the
increments of \((i/k)^r\).  If \(r=0\), a single common tensor suffices.
Thus these fields can be chosen smooth in \(U\), and their child marks
are conditionally independent given the preceding levels.

The perturbation has two tasks.  A term linear in the projected diagonal
overlaps makes those overlaps concentrate, and Gaussian monomial terms
enforce joint GG identities.  We choose their strengths so that both
perturbations are negligible for the pressure, while pressure
concentration still controls the overlap fluctuations.

Set \(e_N=N^{-1/16}\).  For parameters \(u_j,v_a\in[1,2]\), add
\begin{equation}\label{rot:perturbation}
 V_N(\sigma,\alpha)
   =Ne_N\sum_{a=1}^m v_a A^a_{11}(\sigma)
     +\sqrt N e_N\sum_{j\le N}2^{-j}u_jY_j(\sigma,\alpha)
\end{equation}
to the Hamiltonian.  Denote the resulting random pressure by
\(P_N^{\mathrm{enr}}(t,h,u,v)\).  Gibbs brackets include the spin and
cascade label, and \(\E\) averages all disorder.
Direct comparison for the first term and Gaussian comparison for the
second give, uniformly in the displayed parameters,
\begin{equation}\label{rot:small-pressure}
 \left|\E P_N^{\mathrm{enr}}-\E P_N^{\mathrm{enr},0}\right|
 \le C e_N.
\end{equation}
Here the superscript \(0\) denotes omission of \(V_N\), with the same
prior and enrichment.

Conditional on \(U\) and the common data, the spin partition integral
at a leaf is a positive hierarchical mark.  Proposition~\ref{arr:rpc}
therefore applies to this integral, including
\eqref{rot:perturbation}.  In particular,
\begin{equation}\label{rot:level-law}
 \E\langle\ind_{\{c(\alpha^1,\alpha^2)=i\}}\rangle=w_i.
\end{equation}
This is an identity after disorder averaging; no assertion about the
level probabilities in an individual realization is needed.

\subsection{Concentration and overlap identities}

The group Poincar\'e estimate is classical; Haar concentration on the
special orthogonal group is part of the curvature approach of
Gromov--Milman and Bakry--\'Emery \cite{GromovMilman1983,BakryEmery1985}.
We prove the estimate needed here and account for the two components
of the orthogonal group.

\begin{lemma}[Concentration]\label{rot:concentration}
For the preceding models, fix the cascade partition and a bound on
\(h_{k-1}\).  There is a constant independent of \(N\), of the
deterministic prior, and of the remaining parameters in their stated
ranges such that
\begin{equation}\label{rot:variance}
 \Var(P_N^{\mathrm{enr}})\le \frac CN.
\end{equation}
The same assertion holds on any fixed slightly enlarged bounded
interval for the perturbation parameters.  For the pressure of the
original Haar model, without Gaussian enrichment or perturbations,
\begin{equation}\label{rot:fourth}
 \E|P_N-\E P_N|^4\le \frac C{N^2}.
\end{equation}
The constants are independent of deterministic weights in the spin
prior.
\end{lemma}

\begin{proof}
We first record the elementary gradient estimates used below.
For standard Gaussian data \(g\),
\[
 \Var f(g)\le \E|\nabla f(g)|^2.
\]
For completeness, differentiate
\(\E f(g)f(rg+\sqrt{1-r^2}g')\) in \(0<r<1\);
Gaussian integration by parts gives the corresponding scalar product
of gradients, bounded by \(\E|\nabla f(g)|^2\), and integration in \(r\)
proves the inequality.  Approximation extends it from smooth bounded
functions to Lipschitz functions.
Applied to \(f(g/|g|)\), radial independence yields the sphere bound
\[
 \Var_{\mathbb S^{d-1}}f
 \le \frac1{d-2}\E|\nabla_{\mathbb S^{d-1}}f|^2,\qquad d>2.
\]
One may cut off at \(g=0\) before passing to this formula.

Here is a corresponding bound on \(SO(d)\).  Let its Dirichlet form be
the expected sum of squares of the derivatives generated by
column-plane rotations, and let \(C_d\) be a Poincare constant for that
form.  Conditional on a column, the fiber is \(SO(d-1)\).  Its
conditional variance is bounded using the pairs that do not involve
that column.  Moving the column in a unit tangent direction transports
the entire fiber by a rotation; Jensen's inequality bounds the squared
gradient of the conditional mean by the conditional sum of squares of
the derivatives involving that column.  Apply the sphere bound to the
conditional mean and average over the choice of column.  Each pair
appears in \(d-2\) fibers and involves two columns, so
\[
 C_d\le\left(1-\frac2d\right)C_{d-1}
            +\frac2{d(d-2)},\qquad d>2.
\]
The circle has a finite Poincare constant.  Multiplication by
\(d(d-1)\) in this recurrence gives \(C_d=O(d^{-1})\).
The Dirichlet form of a Frobenius-Lipschitz function is bounded by a
fixed multiple of its squared Lipschitz constant.

Conditional on \(U\), apply the backward marked-cascade recursion to
\(Z_N=\exp(NP_N^{\mathrm{enr}})\).  Proposition~\ref{arr:rpc}
gives
\[
 \log Z_N=\mathcal R_N(g_0)+\log S'-\log S,
\]
where \(g_0\) denotes the common Gaussian marks,
\(\mathcal R_N\) is the root recursion value, and \(S,S'\)
are unnormalized cascade totals with the same conditional law.
For the fixed cascade partition, their log variances are bounded,
also conditional on \(U,g_0\).  In particular the conditional
second moment of \(\log S'-\log S\) is bounded uniformly.
The root recursion is
\(C\sqrt N\)-Lipschitz in the common standard Gaussian coefficients.
Indeed each exponent has that coefficient-norm bound: the spin norm
is \(\sqrt N\), the field heights are bounded, and the squared
perturbation coefficients sum to \(O(Ne_N^2)\).  Integration and
log-exponential integration preserve a common Lipschitz bound.
The Gaussian inequality thus bounds the conditional variance of
the normalized log partition by \(C/N\).

It remains to control its conditional mean as a function of \(U\).
For deterministic quadratic terms the normalized exponent is
uniformly Lipschitz in Frobenius distance, since
\(|\sigma|=\sqrt N\).  For the Gaussian terms, compare conditional
covariances.  Each projected overlap is uniformly Lipschitz in \(U\);
the degrees of the monomials grow at most linearly, and
\(\sum_j(1+j)4^{-j}<\infty\).  Gaussian comparison, divided by \(N\),
therefore gives a uniform Lipschitz bound on
\(\E[P_N^{\mathrm{enr}}\mid U]\).  The linear enrichment is independent
of \(U\).  This conditional mean depends only on the spectral
projectors, so multiplication of one spectral row by \(-1\) shows
that it has the same law on both components of \(O(N)\).
The group inequality and conditional variance decomposition prove
\eqref{rot:variance}.

Without the added Gaussian fields, \(P_N\) itself is uniformly
Lipschitz in \(U\) and depends only on those projectors.  Write
\(F=P_N-\E P_N\).  Applying the group inequality first to \(F\) and
then to \(F^2\) gives
\[
 \E F^2\le C/N,\qquad
 \Var(F^2)\le (C/N)\E F^2\le C/N^2.
\]
Since \(\E F^4=\Var(F^2)+(\E F^2)^2\), this proves
\eqref{rot:fourth}.  Smoothing by small group translates justifies
these applications for Lipschitz functions.  All estimates used only
the common spin norm and normalization of the prior.
\end{proof}

We next specify both ways in which the perturbation parameters will
be chosen.  For the upper bound they are deterministic minimizers of
an expected objective.  For the lower bound they are averaged, and we
retain parameter sets on which every subsequential limit has the
required overlap identities.

\begin{proposition}[Perturbations enforcing GG]\label{rot:gg}
Use \eqref{rot:perturbation}, and fix the finite cascade partition and
a compact range of the other parameters.  The following two
procedures give limiting arrays with deterministic diagonals and
joint Ghirlanda--Guerra identities for \((A^1,\ldots,A^m,T)\).
\begin{enumerate}
\item Choose a deterministic minimum of
\[
 -\E P_N^{\mathrm{enr}}(t,h,u,v)+J_N(t,h)
       +\Pi_N(u,v),
 \qquad
 \Pi_N=\sum_{j\le N}2^{-j}(u_j-3/2)^2
             +\sum_a(v_a-3/2)^2,
\]
over any compact set of the other parameters and
\((u,v)\in[1,2]^{N+m}\); here \(J_N\) is independent of \(u,v\).
Every subsequential overlap limit at these minima, after a further
subsequence for the diagonal means, has the stated properties.
\item Fix any deterministic sequence of the other parameters and
integrate over product uniform probability \(\eta\) on
\(\Theta=[1,2]^{\N}\times[1,2]^m\), ignoring \(u_j\) for \(j>N\).
There are measurable sets \(\mathcal G_N\subset\Theta\) with
\(\eta(\mathcal G_N)\to1\) such that every sequence
\(\theta_N\in\mathcal G_N\) has the same properties at each
subsequential overlap limit, after selecting convergent diagonal
means.
\end{enumerate}
These assertions hold for every deterministic prior described above.
They remain valid when \(t=t_N\) and the bounded enrichment parameters
vary with \(N\).
\end{proposition}

\begin{proof}
Let \(g_N=\E P_N^{\mathrm{enr}}\).
For each fixed perturbation coordinate \(d\), convexity of the log
partition makes \(P_N^{\mathrm{enr}}\) and \(g_N\) convex in \(d\).
Expected first derivatives are uniformly bounded on a slightly
enlarged interval.  More precisely, direct differentiation gives
an \(O(e_N)\) bound for a \(v_a\)-derivative, and Gaussian integration
by parts gives \(O_j(e_N^2)\) for a fixed \(u_j\)-derivative.

In the first procedure, comparison with the centered choice of one
coordinate shows that each fixed coordinate lies in the interior
for all large \(N\).  At a minimum, its first derivative matches the
penalty derivative.  Comparison with the choices \(d+s\) and \(d-s\)
therefore gives
\[
 0\le g_N(d+s)-g_N(d)-s g_N'(d)\le C_d s^2
\]
and the analogous negative-increment estimate, whenever the
increments remain in the interval.  In particular
\(0\le g_N''(d)\le C_d\).  Convex secants for the random function,
followed by Lemma~\ref{rot:concentration}, imply
\[
 \E|\partial_dP_N^{\mathrm{enr}}-\partial_dg_N|
 \le C_d\left(s+\frac{N^{-1/2}}s\right).
\]
With \(s=N^{-1/4}\), the right side is \(O_d(N^{-1/4})=o(e_N^2)\).
The parameter minimum is deterministic, so the uniform pointwise
variance estimate applies without conditioning on the disorder.

For the second procedure, integrate the same convex-secant
inequalities in \(d\).  The integral of
\(g_N'(d+s)-g_N'(d-s)\) is \(O_d(s)\), by the endpoint bounds on
the expected first derivatives on the enlarged interval.
Thus
\begin{equation}\label{rot:derivative-fluctuation}
 \int_\Theta
 \E|\partial_dP_N^{\mathrm{enr}}-\partial_dg_N|\,d\eta
 =O_d(N^{-1/4}).
\end{equation}
Convexity and the same endpoint bounds also give
\(\int_\Theta g_N''(d)\,d\eta\le C_d\).

These estimates control both thermal and disorder fluctuations.
For a group coordinate,
\[
 \partial_{v_a}P_N^{\mathrm{enr}}=e_N\langle A^a_{11}\rangle,
 \qquad
 \partial_{v_a}^2g_N
   =Ne_N^2\E\Var_{\langle\cdot\rangle}(A^a_{11}).
\]
Consequently the mean absolute disorder fluctuation of
\(\langle A^a_{11}\rangle\) is \(O(N^{-3/16})\), and its mean
absolute thermal fluctuation is \(O(N^{-7/16})\).
Here and below these bounds hold at a minimum in the first
procedure and after integration in the second.
At the minima, select a convergent subsequence of the deterministic
means \(d_{a,N}=\E\langle A^a_{11}\rangle\).
For the averaged procedure, the means are first kept as functions
\(d_{a,N}(\theta)\); select their limits only after choosing a
sequence \(\theta_N\in\mathcal G_N\) as below.
In either procedure the resulting subsequence satisfies
\begin{equation}\label{rot:fixed-diagonals}
 \E\langle|A^a_{11}-d_a|\rangle\longrightarrow0,
 \qquad d_a\ge0,\quad \sum_a d_a=1.
\end{equation}
For the Gaussian coordinate,
\[
 \partial_{u_j}P_N^{\mathrm{enr}}
     =\frac{e_N2^{-j}}{\sqrt N}\langle Y_j\rangle,
 \qquad
 \partial_{u_j}^2g_N
     =e_N^2 4^{-j}\E\Var_{\langle\cdot\rangle}(Y_j).
\]
Its mean absolute disorder and thermal fluctuations are respectively
\(O_j(N^{5/16})\) and \(O_j(N^{1/16})\).  Since
\(\sqrt N e_N=N^{7/16}\), this proves
\begin{equation}\label{rot:field-fluctuation}
 \E\langle|Y_j-\E\langle Y_j\rangle|\rangle
       =o_j(\sqrt N e_N).
\end{equation}
In the integrated case, the centers in these estimates are always
the expectations at the fixed parameter \(\theta\).
Markov's inequality, followed by a diagonal choice over the finitely
many group coordinates and the countably many fixed indices \(j\),
produces measurable sets \(\mathcal G_N\) on which all these
normalized fluctuation errors tend to zero uniformly.  Their
probabilities tend to one.  This proves the required assertion for
every sequence chosen from the sets, rather than only almost every
fixed parameter.

Finally, let \(\Phi\) be a bounded continuous overlap test on \(n\)
replicas.  Conditional Gaussian integration by parts gives
\begin{align}
 \E\langle Y_j(1)\Phi\rangle
  ={}&\sqrt N e_N2^{-j}u_j
    \E\left\langle\Phi
       \left(\sum_{l=1}^n\kappa_j(1,l)
                         -n\kappa_j(1,n+1)\right)\right\rangle .
 \label{rot:gg-ibp}
\end{align}
Compare this with its one-replica version multiplied by
\(\E\langle\Phi\rangle\), use \eqref{rot:field-fluctuation}, and
then use the deterministic diagonal limits
\eqref{rot:fixed-diagonals}.  Every subsequential limit satisfies
GG for each monomial \(\kappa_j\).
Polynomials are dense in the continuous functions on the compact
overlap state space, so these identities give joint GG for the
entire vector.  The array structure results of
Section~\ref{sec:arrays} now apply.
\end{proof}

All Gaussian differentiations above may first be performed after
restricting the cascade prior to finitely many leaves of total mass
tending to one.  At each \(N\), the variances of the additive
Gaussian fields are bounded.  Jensen's inequality on the normalized
prior bounds moments of inverse Gaussian partition factors, while
Gaussian exponential moments bound the differentiated numerators.
These bounds justify passage from the finite restrictions, also
at tied or zero enrichment levels by one-sided limits.

\subsection{Identities from Haar rotations}

The preceding proposition supplies the deterministic diagonals needed
to close the rotation equations.  We next derive those equations
without any symmetry assumption on the spin prior.

\begin{proposition}[Spectral Ward identities]\label{rot:ward}
Consider a subsequential overlap limit supplied by
Proposition~\ref{rot:gg}, and suppose \(t_N\to t\in[0,1]\).
For distinct group indices \(a,b\), and every bounded continuous
\(\Phi\colon[-1,1]\to\R\),
\begin{align}
 \E\!\left[\Phi(Q_{12})
        (\rho_b A^a_{12}-\rho_a A^b_{12})\right]
 ={}&t(\lambda_a-\lambda_b)
 \E\!\left[\Phi(Q_{12})
        (d_a A^b_{12}+A^a_{12}d_b-2A^a_{13}A^b_{23})\right],
 \label{rot:ward-off}\\
 \rho_b d_a-\rho_a d_b
 ={}&t(\lambda_a-\lambda_b)
       \left(d_a d_b-\E[A^a_{12}A^b_{12}]\right).
 \label{rot:ward-diag}
\end{align}
The assertion includes equal eigenvalues in distinct groups.
\end{proposition}

\begin{proof}
Choose row indices \(i\in I_{a,N}\), \(j\in I_{b,N}\), and
differentiate under the plane rotation
\(x_i'=x_j,\ x_j'=-x_i\).
Apply this derivative to
\(\langle x_i^1x_j^2\Phi(Q_{12})\rangle\), keeping the Gaussian
coefficient tensors and the independent linear enrichment fixed.
The Haar-averaged derivative is zero.  The total overlap \(Q_{12}\)
is rotation invariant, so no derivative hits \(\Phi\).
Summing the coordinate derivative over \(i,j\) and dividing by \(N^2\)
gives
\[
 \rho_{a,N}A^b_{12}-\rho_{b,N}A^a_{12}.
\]
The derivative of the main Hamiltonian is
\(t_N(\lambda_a-\lambda_b)x_i x_j\).
The two-replica Gibbs-ratio derivative thus contributes
\[
 t_N(\lambda_a-\lambda_b)
  \left(A^a_{11}A^b_{12}
        +A^a_{12}A^b_{22}-2A^a_{13}A^b_{23}\right).
\]

We check that the remaining terms disappear at the required scale.
The deterministic perturbation derivative has coefficient
\(2e_N(v_a-v_b)\) multiplying \(x_i x_j\), and its normalized
contribution is \(O(e_N)\).
For a Gaussian perturbation derivative, integrate by parts in its
coefficient tensor, conditionally on \(U\) and the cascade.
Pairing with the Gibbs exponent supplies coefficient
\(Ne_N^2 4^{-j}u_j^2\).
Each one-endpoint derivative of its covariance is a sum of at most
\(C(1+j)\) terms consisting of bounded overlap factors times
\[
 \frac{x_j^u x_i^v}{N}
       \quad\hbox{or}\quad \frac{x_i^u x_j^v}{N}.
\]
Multiplication by \(x_i^1x_j^2\), summation over \(i,j\), and
division by \(N^2\) contract these coordinates into products of
normalized projected inner products, whose absolute values are at
most one.  The total error is bounded by
\[
 Ce_N^2\sum_{j\ge1}(1+j)4^{-j}=O(e_N^2).
\]
The independent linear enrichment contributes nothing.
Finite cascade restrictions justify the differentiations as above.
Pass to the limit and use the deterministic diagonals to obtain
\eqref{rot:ward-off}.

The same calculation with \(x_i^1x_j^1\) and a one-replica Gibbs
integral gives
\eqref{rot:ward-diag}.  In this calculation the exponent derivative
is still \(t_N(\lambda_a-\lambda_b)x_i x_j\), and the Gibbs derivative
has one subtraction replica.  This accounts for the coefficient
without an additional factor of two.
\end{proof}

\subsection{Products of replica paths}

The product and symbol below are familiar from Parisi-matrix algebra and the
replica Fourier transform; see M\'ezard--Parisi
\cite[Appendix~II]{MezardParisi1991} and M\"uller--Pankov
\cite[Appendix~E]{MullerPankov2007}.
We derive the corresponding product directly from synchronized overlap
identities, including atoms and jumps.  Our immediate aim is to regard
the diagonal and conditional off-diagonal expressions in the Ward
identities as one product of paths, and then turn that product into
ordinary multiplication by taking symbols.

By joint GG and synchronization, the group off-diagonal overlaps
have simultaneous nondecreasing quantiles \(p_a\), with
\(\sum_a p_a=p\).  In particular each \(p_a\) is constant on every
plateau of \(p\).  We encode a diagonal and an off-diagonal path by
\(L=(d,l)\), and define its symbol by
\begin{equation}\label{rot:symbol}
 \widehat L(s)=d-sl(s)-\int_s^1 l(u)\,du,\qquad 0<s<1.
\end{equation}
We choose right-continuous versions where appropriate.  The symbol alone
does not determine a path: replacing \((d,l)\) by \((d+c,l+c)\)
leaves it unchanged.  We therefore retain the root value \(l(0+)\)
as well.  For the total overlap,
\[
 \widehat Q(s)=1-sp(s)-\int_s^1p(u)\,du=D_p(p(s)).
\]

\begin{lemma}[Replica product]\label{rot:product}
For the synchronized paths \(L_a=(d_a,p_a)\) and \(L_b=(d_b,p_b)\),
the diagonal and conditional off-diagonal expressions on the right
of \eqref{rot:ward-diag} and \eqref{rot:ward-off}, respectively,
form a path \(L_a\star L_b\) with diagonal
\[
 d_a d_b-\int_0^1p_a(u)p_b(u)\,du
\]
and off-diagonal
\begin{equation}\label{rot:product-path}
 \widehat L_a(s)p_b(s)+p_a(s)\widehat L_b(s)
       +sp_a(s)p_b(s)-\int_0^s p_a(u)p_b(u)\,du.
\end{equation}
Its symbol and root are
\begin{align*}
 \widehat{L_a\star L_b}&=\widehat L_a\,\widehat L_b,\\
 (L_a\star L_b)(0+)
   &=\widehat L_a(0+)p_b(0+)+p_a(0+)\widehat L_b(0+).
\end{align*}
\end{lemma}

\begin{proof}
Conditional on \(Q_{12}=q\), GG gives the marginal of \(Q_{13}\)
as one half of the law of \(Q_{12}\) plus one half of \(\delta_q\).
Ultrametricity forces \(Q_{13}=Q_{23}\) below \(q\), and forces the
other overlap to equal \(q\) when one exceeds \(q\).
To include atoms explicitly, suppose the quantile equals \(q\) on
\((\alpha,\beta)\).  Write \(a_*\) and \(b_*\) for the constant
values of \(p_a\) and \(p_b\) on that interval.  The preceding
conditional laws give
\[
 \begin{split}
 2\E[A^a_{13}A^b_{23}\mid Q_{12}=q]
 ={}&\int_0^\alpha p_a(u)p_b(u)\,du
       +b_*\int_\beta^1p_a(u)\,du
       +a_*\int_\beta^1p_b(u)\,du\\
    &+(2\beta-\alpha)a_*b_* .
 \end{split}
\]
This is exactly the expression subtracted from
\(d_a p_b+p_a d_b\) in \eqref{rot:product-path}, evaluated as
\(s\uparrow\beta\).  Formula \eqref{rot:product-path} is constant across
the plateau.  The same disintegration with \(\alpha=\beta\)
gives the formula at non-atomic values, for almost every
conditioning value.

For the algebraic claims, the Stieltjes identity
\[
 d\widehat L=-s\,dl
\]
follows from \eqref{rot:symbol}.  If \(c\) denotes
\eqref{rot:product-path}, its jump at a common discontinuity \(v\)
is
\[
 \Delta c=\widehat L_a(v-)\Delta p_b
                  +(\Delta p_a)\widehat L_b(v+).
\]
Its continuous Stieltjes part obeys the corresponding product rule.
Thus \(d\widehat{L_a\star L_b}
=d(\widehat L_a\widehat L_b)\).
At \(1-\), subtracting the off-diagonal from the diagonal gives
\((d_a-p_a(1-))(d_b-p_b(1-))\), the required endpoint product.
The symbols consequently agree everywhere up to the choice of
versions.  Taking \(s\downarrow0\) in \eqref{rot:product-path}
gives the root formula.
\end{proof}

\subsection{Reconstruction in each spectral group}

Let \(b_t(x)\) be the inverse resolvent for the finite law
\(\sum_a\rho_a\delta_{t\lambda_a}\), and set
\begin{equation}\label{rot:f}
 f_a(x)=\frac{\rho_a}{b_t(x)-t\lambda_a}\quad(x>0),
 \qquad f_a(0)=0.
\end{equation}
For a finite alphabet of positive proportions,
\(b_t(x)>\max_a t\lambda_a\) for \(x>0\).  The functions \(f_a\)
are smooth jointly in \(x,t\), including \(x=0\): near zero, apply
the inverse function theorem with \(1/b_t\) as the inverse variable.
Moreover
\begin{equation}\label{rot:f-derivative}
 f_a'(x)=
 \frac{\rho_a/(b_t(x)-t\lambda_a)^2}
      {\sum_b\rho_b/(b_t(x)-t\lambda_b)^2},
 \qquad
 f_a'(0)=\rho_a,\qquad \sum_a f_a'(x)=1 .
\end{equation}

\begin{proposition}[Projected-overlap paths]\label{rot:paths}
Under the hypotheses of Proposition~\ref{rot:ward}, let \(p\)
be the quantile of the total overlap.  Then the entire group path
is determined by \(p\):
\begin{equation}\label{rot:full-paths}
 p_a(s)=\int_0^{p(s)}f_a'(D_p(r))\,dr,\qquad
 d_a=\int_0^1f_a'(D_p(r))\,dr .
\end{equation}
These formulas hold at \(t=0\), for repeated eigenvalues, and when
\(D_p(p(s))=0\).  More generally, for a smooth scalar or matrix
function \(f\) on \([0,1]\), the path
\begin{equation}\label{rot:function-path}
 l_f(s)=\int_0^{p(s)}f'(D_p(r))\,dr,\qquad
 d_f=f(0)+\int_0^1f'(D_p(r))\,dr
\end{equation}
has symbol \(f(D_p(p(s)))\) and root
\(p(0+)f'(D_p(p(0+)))\).  Symbol and root uniquely determine
a bounded-variation path.
\end{proposition}

\begin{proof}
By Lemma~\ref{rot:product}, the two Ward identities combine into
\[
 \rho_bL_a-\rho_aL_b
      =t(\lambda_a-\lambda_b)(L_a\star L_b).
\]
Set \(\alpha_a=\widehat L_a(s)\) and \(x=\sum_a\alpha_a\).
GG positivity and Gram bounds give \(0\le p_a(s)\le d_a\),
so \(\alpha_a\ge0\),
and \(x=D_p(p(s))\).
If \(x>0\), every \(\alpha_a\) is positive: a zero component
paired with a positive component would give
\(-\rho_a\alpha_b=0\).
Dividing the symbol equations therefore gives a common number
\[
 \frac{\rho_a}{\alpha_a}+t\lambda_a=b
       \quad\hbox{for every }a.
\]
The positive sum constraint is the resolvent equation for \(b_t(x)\).
Hence \(\alpha_a=f_a(x)\).  If \(x=0\), all the symbols are zero,
which is the same conclusion.

At the root, put \(r_a=p_a(0+)\), \(p_0=\sum_a r_a\), and
\(x=D_p(p_0)\).  The root equations are
\[
 \rho_b r_a-\rho_a r_b
 =t(\lambda_a-\lambda_b)
        \bigl(f_a(x)r_b+r_a f_b(x)\bigr).
\]
For \(x>0\), substitution of \eqref{rot:f} shows that
\[
 \frac{(b_t(x)-t\lambda_a)^2}{\rho_a}\,r_a
\]
is independent of \(a\).  The sum constraint and
\eqref{rot:f-derivative} give \(r_a=p_0f_a'(x)\).
For \(x=0\), the root equations directly give
\(\rho_b r_a=\rho_a r_b\), and therefore
\(r_a=\rho_a p_0=p_0 f_a'(0)\).

It remains to recover the paths, including their diagonals, from
these data.  Let \(F_p(r)=|\{s:p(s)\le r\}|\).  Then
\(D_p'(r)=-F_p(r)\) almost everywhere and \(D_p(1)=0\).
For the path in \eqref{rot:function-path}, Fubini's theorem gives
\[
 \begin{split}
 \widehat L_f(s)
 &=f(0)+\int_{p(s)}^1 F_p(r)f'(D_p(r))\,dr\\
 &=f(D_p(p(s))).
 \end{split}
\]
Since \(D_p\) is constant on \([0,p_0]\), its root is
\(p_0 f'(D_p(p_0))\).
Finally, equal symbols imply equal Stieltjes derivatives of the
off-diagonal paths on every \([\epsilon,1)\), because
\(d\widehat L=-s\,dl\).  The paths differ by a constant, fixed by
the root, and the symbol then fixes the diagonal.  This also treats
jumps accumulating at zero.  Applying this calculation to \(f_a\)
proves \eqref{rot:full-paths}.

No step divided by an eigenvalue difference.  At \(t=0\),
\(f_a(x)=\rho_a x\), so \(p_a=\rho_a p\) and \(d_a=\rho_a\).
Tied eigenvalues similarly give proportional paths in the tied
groups.  If \(p\equiv1\), then \(D_p\equiv0\) and
\(p_a=d_a=\rho_a\), also covered by the formulas.
\end{proof}

For later Gaussian marks, the diagonal excess in group \(a\) is
\[
 d_a-p_a(1-)=\int_{p(1-)}^1 f_a'(D_p(r))\,dr\ge0.
\]
After division by \(\rho_a\), the corresponding function in
\eqref{rot:function-path} is
\(H_a(x)=(b_t(x)-t\lambda_a)^{-1}\), with \(H_a(0)=0\).
Its derivative is bounded by \(1/\rho_a\).
Thus the group formulas give uniformly bounded Gaussian
covariances for all trial paths, with zero residual variance when
the top overlap is one.
